\documentclass[10pt]{article}

\usepackage[utf8]{inputenc}

\usepackage[T1]{fontenc}

\usepackage{amsmath}

\usepackage{amsfonts}

\usepackage{amssymb}

\usepackage[version=4]{mhchem}

\usepackage{stmaryrd}

\usepackage{mathrsfs}

\usepackage{bbold}

\usepackage{graphicx}

\usepackage[export]{adjustbox}

\graphicspath{ {./images/} }



\begin{document}

трич. пространство, ва $X$ возникает равномерность, порожденная метрикой $\rho$. Систему окружений этой равномерности образуют всевозможные множества, содержащие муожества вида $\{(x, y): \rho(x, y)<\varepsilon\}, \varepsilon>0$. При әтом топологии на $X$, индулированные метрикой и равномерностью, совпадают. Равномерные структуры, порожденные метриками, наз. ме т р и з у е м ы ми. Р. п. были введены в 1937 А. Вейлем [1] (посредством окружений; определение Р. п. посредством равномерных покрытий было дано в 1940, см. [4]). Однако идея использования многократной звездной вписанности для построения функций появилась ранее у Л. С. Понтрягина (см. [5]) (впоследствии эта идея была использована при доказательстве полной регулярности топологии отделимого Р..). Первоначально равномерные структуры использовались как инструмент для изучевия (порожденных ими) топологий (подобно тому, как метрика на метризуемом пространстве часто используется для изучения топологич. свойств этого пространства). Однако теория Р. П. имеет и самостоятельное значение, хотя и тесно связана с теорией топологич. пространств.



Отображение $f: X \rightarrow Y$ Р. П. $X$ в Р. П. $Y$ наз. р а вно ме р н о не і р е ры в ным, если для любого равномерного покрытия $\alpha$ пространства $Y$ система $f^{-1} \alpha=$ $=\left\{f^{-1} U: U \in \alpha\right\}$ являөтся равномерным покрытием $X$. Всякое равномерно непрерывное отображение является непрерывным относительно топологий, порожденных равномерными структурами на $X$ и $Y$. Если равномерные структуры на $X$ и $Y$ индуцированы метриками, то равномерно непрерывное отображение $f: X \rightarrow Y$ оказывается равномерно нешрерывным в классич. смысле как отображение метрич. пространств.



Более содержательной является теория Р. п., к-рая удовлетворяет дополнительной аксиоме отделимости:



U4) $\bigcap_{V \in \mathscr{N}} V=\Delta$



(в терминах окружений) или (в терминах равномерных покрытий):



С3) для любых двух точек $x, y \in X, x \neq y$, существует такое $\alpha \in(5$, что викакой алемент $\alpha$ не содержит точки $x$ и $у$ одновременно.



Далее речь будет идти только о Р. п., наделенных отделимой равномерной структурой. Топология, порожденная на $X$ отделимой равномерностью, является вполне регулярной и обратно, всякая вполне регулярная топология на $X$ порождается нек-рой отделимой равномерной структурой. Как правило, существует много различных равномерностей, порождающих одинаковую топологию на $X$. В частности, метризуемая топология может порождаться неметризуемой отделимой равномерностыо.



Р. п. $(X, \mathfrak{A})$ является метризуемым тогда и только тогда, когда $\mathfrak{Y}$ имеют счетную базу. При этом б а 3 о й равномерности наз. (в терминах окружений) всякая подсистөма $\mathfrak{B} \subset \mathfrak{A}$, удовлетворяющая условию: для любого $V \in \mathfrak{A}$ существует такое $W \in \mathfrak{B}$, что $W \subset V$, или (в терминах равномерных покрытий) подсистема $\mathfrak{X} \subset \mathbb{S}$ такая, что для любого $\alpha \in \mathbb{S}$ существует $\beta \in \mathfrak{A}$, к-рое вписано в $\alpha$. В е с о м P. п. $(X, \mathfrak{U})$ наз. наименьшая мощность базы равномерности $\mathfrak{A}$.



\includegraphics[max width=\textwidth, center]{2023_07_08_6945767785c131642076g-1}

жений $\mathfrak{A}_{M}=\{(M \times M) \cap V: V \in \mathfrak{A}\}$ определяет равномерность на $M$. Пара $\left(M, \mathfrak{Y}_{M}\right)$ наз. п о д п р о с т р а н с тв о м P. п. $(X, \mathfrak{Y})$. Отображение $f: X \rightarrow Y$ Р. п. $(X, \mathfrak{X})$ в Р. п. $\left(Y, \mathfrak{U}^{\prime}\right)$ наз. р а в н о м е р ны м в ло ж ен и е м, если $f$ взаимно однозначно, равномерно непрерывно и отображение $f^{-1}:\left(f X, \mathfrak{U}_{f X}^{\prime}\right) \rightarrow(X, \mathfrak{U})$ также равномерно непрерывно.



Р. п. $X$ наз. п о л н ы м, если всякий фильтр Коши в $X$ (т. е. фильтр, содержащий нек-рый әлемент всякого равномерчого покрытия) имеет точку прикосновения (т. е. точку, лежащую в пересечении замыканий әлементов фильтра). Метризуемое Р. п. является полным тогда и только тогда, когда полна метрика, порождающая его равномерность. Јюбое Р. п. $(X, \mathfrak{U})$ может быть равномерно вложено в качестве всоду плотного подмножества в единственное (с точностью до равномерного изоморфизма) полное Р. п. $(\tilde{X}, \widehat{\mathfrak{Y}})$, к-рое наз. п оп о л п е н и е м $(X, \mathfrak{A})$. Топология пополнения $(\tilde{X}$, $\widetilde{\mathfrak{H}})$ P. II. $(X, \mathfrak{A})$ бикомпактна тогда и только тогда, когда $\mathfrak{U}$ является п р е ко м п а к т н о й р а в ном е р н о с т ь ю (т. е. такой, что в любое равномерное покрытие можно вписать конечное равномерное покрытие). В этом случае пространство $\tilde{X}$ является бикомпактным расширением $X$ и наз. р а с пा и р е н и е м С а м ю э л я пространства $X$ относительно равномерности $\mathfrak{U}$. Для всякого бикомпактного расширения $b X$ пространства $X$ существует единственная прекомпактная равномерность на $X$, расширение Самюэля относительно к-рой совпадает с $b X$. Таким образом, на языке прекомпактных равномерностей описываются все бикомпактные расширения $b X$. На бикомпактном пространстве существует единственная равномерность (полная и прекомпактная).



Всякая равномерность $\mathfrak{X}$ на множестве $X$ индуцирует близость $\delta$ по следующей формуле:



\[

A \delta B \Leftrightarrow(A \times B) \cap V \neq \varnothing

\]



для любого $V \in \mathfrak{A}$. При этом топологии, порожденные на $X$ равномерностью $\mathfrak{U}$ и близостью $\delta$, совпадают. Любое равномерно непрерывное отображение является близостно нешрерывным относительно близостей, порожденных равномерностями. Как правило, существует много различных равномерностей, порождающих на $X$ одну и ту же близость. Тем самым множество равномерностей на $X$ распадается на классы эквивалентности (две равномерности эквивалентны, если индуцированные ими близости совпадают). Каждый класс әквивалентных равномерностей содергит ровно одну прекомпактную равномерность, причем расширения Самюэля относительно этой равномерности совпадают с расширениями Смирнова (см. Близости пространство) относительно близости, индуцированной равномерностями этого класса. В множестве равномерностей на $X$ задается естественный частичный порядок: $\mathfrak{A}>\mathfrak{Y} \mathfrak{Y}^{\prime}$, если $\mathfrak{A} \supset \mathfrak{A} \mathfrak{Y}^{\prime}$. Среди всех равномерностей, порождающих на $X$ фиксированную топологию, есть наибольшая - т. н. у н ив е р с а льва я ра в оме р ность, к-рая индуцирует на $X$ близость Стоуна - Чеха. Всякая прекомпактная равномерность является наименьшим элементом в классе эквивалентных ей равномерностей. Если (5 - система равномерных покрытий нек-рой равномерности на $X$, то система равномерных покрытий әквивалентной прекомпактной равномерности состоит из таких покрытий $X$, в к-рые можно вписать конечные покрытия из (5.



П р о и з в еде н и е м Р. п. $\left(X_{t}, \mathfrak{A}_{t}\right), t \in T$, наз. Р. п. $\left(\Pi X_{t}, \Pi \mathfrak{Y} \mathfrak{Y}_{t}\right)$, где $\Pi \mathfrak{Y}_{t}-$ равномерность на $X_{t}$, базу окружений к-рой образуют множества вида



\[

\begin{gathered}

\left\{\left(\left\{x_{t}\right\},\left\{y_{t}\right\}\right):\left(x_{t_{i}}, y_{t_{i}}\right) \in V_{t_{i}}, i=1, \ldots, n\right\}, \\

t_{i} \in T, V_{t_{i}} \in \mathfrak{A}_{t_{i}}, n=1,2, \ldots .

\end{gathered}

\]



Топология, индуцированная на $\Pi X_{t}$ равномерностью $\Pi \mathfrak{U}_{t}$, совпадает с топологией тихоновского произведения пространств $X_{t}$. Проекции произведения $\mathrm{P}$. п. на сомножители равномерно непрерывны. Всякое Р. П. веса $\tau$ может быть вложено в произведение $\tau$ экземпляров метризуемых Р. П.



Семейство $F$ непрерывных отображений топологич. пространства $X$ в P. П. $(Y, \mathfrak{U})$ наз. р а в н о с т е пе н-





\end{document}